\section{The actual coordinate sections give canonical original projective line frames}

\begin{theorem}[The same original coordinate section is a true frame on its chart]
\label{thm:fc-coordinate-section-own-chart-frame}
\leanok
\lean{EnvelopeProperty.ProjectiveCoordinateSectionRelations.sectionHom_isIso_on_own_chart}
\uses{thm:fc-generated-tuple-twist-section-scalars}
The actual original global section $X_i$ of the SAME original
$O(1)$ is a whole line isomorphism after native restriction to its
actual original chart $D_+(X_i)$, for every original coefficient
ring and every original choice of the coordinate hyperplane model
of $O(1)$.
\end{theorem}
\begin{proof}
Its genuine original standard-frame coefficient is the actual
fraction $X_i/X_i=1$. The original local scalar identity therefore
makes the section followed by the genuine original frame the
identity. The frame is already a genuine isomorphism, so that
SAME original section is an isomorphism too.
\end{proof}

\begin{theorem}[All same original coordinates retain their whole ratio relation]
\label{thm:fc-coordinate-section-chart-ratio}
\leanok
\lean{EnvelopeProperty.ProjectiveCoordinateSectionRelations.sectionHom_ratio_on_chart}
\uses{thm:fc-coordinate-section-own-chart-frame}
On the SAME original coordinate chart, the original whole section
$X_j$ is the original whole section $X_i$ multiplied by its actual
original section-ring ratio $X_j/X_i$, as an entire native
module-morphism identity.
\end{theorem}
\begin{proof}
Both sections have their genuine original coefficients in the SAME
actual line frame. The first coefficient is $X_j/X_i$ and the
pivot coefficient is one. After composing with that genuine frame,
the asserted original morphism identity is scalar multiplication
by $X_j/X_i$ on the structure line. Cancel the actual original
frame to get the whole original section identity.
\end{proof}
